\begin{lemma}
\label{lemma-non-descending-property-projective}
The properties
\begin{enumerate}
\item[] $\mathcal{P}(f) =$``$f$ is projective'', and
\item[] $\mathcal{P}(f) =$``$f$ is quasi-projective''
\end{enumerate}
are not Zariski local on the base. A fortiori, they are not fpqc local
on the base.
\end{lemma}

\begin{proof}
Following Hironaka \cite[Example B.3.4.1]{H},
we define a proper morphism of smooth complex 3-folds
$f:V_Y\to Y$
which is Zariski-locally projective, but not projective. Since $f$ is proper
and not projective, it is also not quasi-projective.

\medskip\noindent
Let $Y$ be projective 3-space over the complex numbers ${\mathbf C}$.
Let $C$ and $D$ be smooth conics in $Y$ such
that the closed subscheme $C\cap D$ is reduced and consists
of two complex points $P$ and $Q$. (For example,
let $C=\{ [x,y,z,w]: xy=z^2, w=0\}$, $D=\{ [x,y,z,w]:
xy=w^2, z=0\}$, $P=[1,0,0,0]$,
and $Q=[0,1,0,0]$.) 
On $Y-Q$, first blow up the curve $C$, and then blow
up the strict transform of the curve $D$ (Divisors, Definition 
\ref{divisors-definition-strict-transform}). On $Y-P$, first blow up
the curve $D$, and then blow up the strict transform of the curve
$C$. Over $Y-P-Q$, the two varieties we have constructed are canonically
isomorphic, and so we can glue them over $Y-P-Q$. The result
is a smooth proper 3-fold $V_Y$ over ${\mathbf C}$. The morphism
$f:V_Y\to Y$ is proper and Zariski-locally projective (since
it is a blowup over $Y-P$ and over $Y-Q$), by Divisors,
Lemma \ref{divisors-lemma-blowing-up-projective}. We will show that
$V_Y$ is not projective over ${\mathbf C}$. That will imply that
$f$ is not projective.

\medskip\noindent
To do this, let $L$ be the inverse image in $V_Y$ of a complex point
of $C-P-Q$, and $M$ the inverse image of a complex point of $D-P-Q$.
Then $L$ and $M$ are isomorphic to the projective line
${\mathbf P}^1_{{\mathbf C}}$. 
Next, let $E$ be the inverse image in $V_Y$ of $C\cup D\subset Y$ in $V_Y$;
thus $E\rightarrow C\cup D$ is a proper morphism, with fibers
isomorphic to ${\mathbf P}^1$ over $(C\cup D)-\{P,Q\}$.
The inverse
image of $P$ in $E$ is a union of two lines $L_0$ and $M_0$, and we have
rational equivalences of cycles $L\sim L_0+M_0$ and $M\sim M_0$ on $E$
(using that $C$ and $D$ are isomorphic to ${\mathbf P}^1$).
Note the asymmetry resulting from the order in which we blew
up the two curves. Near $Q$, the opposite happens. So the inverse image
of $Q$ is the union of two lines $L_0'$ and $M_0'$, and we have
rational equivalences $L\sim L_0'$ and $M\sim L_0'+M_0'$ on $E$.
Combining these equivalences, we find that $L_0+M_0'\sim 0$
on $E$ and hence on $V_Y$. If $V_Y$ were projective over ${\mathbf C}$,
it would have
an ample line bundle $H$, which would have degree $> 0$ on all curves
in $V_Y$. In particular $H$ would have positive degree on $L_0+M_0'$,
contradicting that the degree of a line bundle is well-defined
on 1-cycles modulo rational equivalence on a proper scheme
over a field (Chow Homology,
Lemma \ref{chow-lemma-proper-pushforward-rational-equivalence}
and Lemma \ref{chow-lemma-factors}).
So $V_Y$ is not projective over ${\mathbf C}$.
\end{proof}
